\chapter{虚拟内存与地址空间}

\section{虚拟地址为何需要独立解释}

先看一个容易误判的程序：

\begin{lstlisting}
char* p = malloc(1 << 30);   // request 1 GiB virtual range
p[0] = 1;                    // touch first page
p[4096] = 2;                 // touch second page
\end{lstlisting}

如果 \code{malloc} 成功就必须立即占用 1 GiB 连续物理内存，那么程序只触碰两个页面也会
长期占住其余空间；多个进程还会因为各自期望使用相似地址而互相冲突。系统需要先给程序
一个可使用的地址范围，在真正访问其中某页时再决定该页对应哪一个物理页。与此同时，
程序不能借此访问内核或其他进程，系统还必须为每个范围保存读、写和执行权限。

由这两个要求才得到\emph{虚拟内存}：进程使用的是自身地址空间中的虚拟地址，页表在需要时
把虚拟页映射到物理页，并附带 readable、writable、executable、present 等状态。一次访问
是否合法，不只取决于地址数值，也取决于该范围是否已经获得承诺、访问类型是否被允许，
以及尚未准备物理页时能否通过缺页处理补齐。

把这条路径拆开，会看到三个层次：

\begin{enumerate}
\tightlist
\item 用户代码只拿到虚拟地址，例如 \code{p + 4096}。
\item 内核用 VMA 记录“这段地址理论上属于堆，允许读写，可以按需分配”。
\item CPU/MMU 用页表项记录“这一页当前是否有物理页，权限是什么”。
\end{enumerate}

因此“没有 PTE”和“非法地址”不是同一件事。一个地址如果不在任何 VMA 中，访问就是非法；一个地址在 VMA 中但还没有 PTE，可能只是第一次访问，需要缺页处理分配或调入页面。理解这层区别，后面的 \code{mmap}、COW、swap、page cache 才不会混在一起。

隔离来自默认不可见。一个进程不能随便读写另一个进程的页；没有映射的地址会触发缺页；只读页不能写；不可执行页不能取指。操作系统通过这些规则同时提供保护、共享和按需分配。

虚拟地址空间不仅是页表。内核还维护 VMA 或等价区域结构，用来描述一段虚拟地址的来源和权限：匿名内存、栈、堆、共享库、文件映射、设备映射。页表项是某一页当前是否能翻译，VMA 是这一段地址理论上是否合法、缺页时应该怎样处理。没有页表项不一定是错误，可能只是尚未按需分配。

一次访问的大致路径是：CPU 发出虚拟地址；TLB 命中则直接得到物理地址和权限；TLB 未命中则硬件或内核进行页表遍历；页表项不存在或权限不符时触发异常；内核根据 VMA 判断是合法缺页、权限错误还是非法访问；合法缺页可能分配物理页、从文件读页、建立 COW 私有页或映射零页。

这条路径可以用一个具体地址走一遍。假设页大小是 4096，进程访问 \code{0x401234}：

\begin{lstlisting}
virtual address = 0x401234
virtual page    = 0x401
offset in page  = 0x234

find VMA containing 0x401234
  yes: [0x400000, 0x500000) readable executable file mapping

look up PTE for virtual page 0x401
  present: physical page 0x83000, readable executable

physical address = 0x83000 * 4096 + 0x234
\end{lstlisting}

若 PTE 不存在，内核不会立刻把程序杀掉，而是先问 VMA：这是不是文件映射的合法页？是不是栈增长？是不是匿名页按需分配？若 VMA 也找不到，才是典型的非法访问。若 VMA 找到了但访问类型不允许，例如往只读代码页写，就是权限错误。缺页异常只是入口，真正结论由 VMA、PTE、访问类型和错误码共同决定。

copy-on-write 是虚拟内存最有代表性的协议。\code{fork} 时，父子进程可以先共享物理页，并把页表标记为只读和 COW。任何一方写入时触发保护异常，内核复制物理页，更新写入方页表，再继续执行。这样 fork 不需要立刻复制整个地址空间，但必须正确维护引用计数和权限。

\begin{keyidea}
虚拟内存不是“把地址翻译成另一个地址”这么简单。它同时维护三件事：哪些虚拟区间被承诺，当前哪些页已经物化成物理页，每次读写执行是否被授权。
\end{keyidea}

\section{VMA、页表项与物理页的职责}

实验的 \code{PageTable} 支持按页映射。读访问要求 present 和 readable；写访问要求 writable；执行访问要求 executable。地址翻译保留页内 offset，帮助读者理解“页号转换，页内偏移不变”。

\begin{longtable}{p{0.25\linewidth}p{0.55\linewidth}}
\toprule
输入 & 计算 \\
\midrule
page size & 4096 \\
virtual address & \code{0x401234} \\
virtual page number & \code{0x401} \\
offset & \code{0x234} \\
physical page number & 查页表得到 \\
physical address & \code{physical\_page * 4096 + offset} \\
\bottomrule
\end{longtable}

权限矩阵要单独写：

\begin{longtable}{p{0.13\linewidth}p{0.14\linewidth}p{0.15\linewidth}p{0.15\linewidth}p{0.15\linewidth}p{0.16\linewidth}}
\toprule
访问 & present & readable & writable & executable & 结果 \\
\midrule
load & yes & yes & no & no & 成功 \\
store & yes & yes & no & no & 权限错误 \\
fetch & yes & yes & no & no & NX 错误 \\
fetch & yes & yes & no & yes & 成功 \\
load & no & - & - & - & 缺页或非法 \\
\bottomrule
\end{longtable}

观察真实系统时，要区分虚拟大小、RSS、page cache、minor fault 和 major fault。虚拟地址空间很大，不代表物理内存已经占用；RSS 增加，可能是匿名页，也可能是文件页；minor fault 可能只是建立映射，major fault 更可能等待存储设备。

\section{页面建立、复制、回收与换出}

\subsection{物理页分配：free list、bitmap 和 buddy}

最简单的物理页分配器是 free list。每个空闲页框放进链表，分配时弹出一个，释放时放回去。它的分配和释放都是 \code{O(1)}，适合固定大小页框；缺点是难以快速找到连续页，也难以检测重复释放和碎片分布。

\begin{lstlisting}
page_alloc():
  if free_list.empty(): return null
  page = free_list.pop()
  page.refcount = 1
  return page

page_free(page):
  assert(page.refcount == 0)
  free_list.push(page)
\end{lstlisting}

bitmap 分配器用一个 bit 表示一个页框是否空闲。它节省元数据，容易扫描连续区域，但朴素扫描是 \code{O(n)}。可以通过分层 bitmap 或缓存上次搜索位置降低平均成本。

buddy allocator 用 2 的幂大小管理连续页块。分配 8 页时，寻找 order=3 的块；如果只有更大的块，就不断拆分；释放时，如果相邻 buddy 也空闲，就合并成更大块。buddy 的分配和释放大约是 \code{O(log N)}，能较好处理连续页需求，但仍会产生内部碎片。

\subsection{虚拟地址查找：VMA 和页表}

页表回答“这一页当前映射到哪里”，VMA 回答“这一段地址理论上是否合法”。缺页时，内核先查 VMA：若地址不属于任何 VMA，就是非法访问；若属于可按需分配的匿名 VMA，就分配零页；若属于文件映射 VMA，就从 page cache 找页或发起 I/O。

VMA 可以用有序区间结构保存。查找地址属于哪个区间，朴素线性扫描是 \code{O(n)}；平衡树查找是 \code{O(log n)}；现代内核还会使用更适合区间查询和缓存局部性的结构。教学模型可以先用有序数组，因为它能把语义讲清楚。

\begin{lstlisting}
handle_page_fault(addr, access):
  vma = find_vma(addr)
  if vma == null: signal(SIGSEGV)
  if !vma.allows(access): signal(SIGSEGV)
  page = materialize_page(vma, addr)
  install_pte(addr, page, vma.permissions)
\end{lstlisting}

\subsection{COW 的引用计数}

COW 的不变量是：共享页必须只读，物理页有引用计数；写入共享页时，写入方得到新副本。若 refcount 为 1，可以直接把页改成可写；若 refcount 大于 1，必须复制。

\begin{lstlisting}
handle_write_fault(vaddr):
  pte = lookup_pte(vaddr)
  page = pte.page
  if !pte.cow: signal(SIGSEGV)
  if page.refcount == 1:
    pte.writable = true
    pte.cow = false
  else:
    new_page = alloc_page()
    copy_page(new_page, page)
    page.refcount -= 1
    install_pte(vaddr, new_page, writable=true, cow=false)
\end{lstlisting}

这段算法最容易错在引用计数和权限更新顺序。若忘记把共享页设为只读，写入不会触发 fault，父子进程会互相污染；若复制后忘记减少旧页引用计数，会泄漏；若 TLB 中仍缓存旧权限，必须刷新或失效对应条目。

用一个父子进程的实际页来走一遍。父进程有一页虚拟地址 \code{0x7000}，物理页是 \code{P42}，内容是字符串 \code{"abc"}。\code{fork} 前，父进程 PTE 指向 \code{P42}，权限可读可写，\code{P42.refcount = 1}。\code{fork} 时，内核不复制 \code{P42} 的 4096 字节，而是让父子页表都指向同一个 \code{P42}，同时把两个 PTE 都改成只读，并在软件元数据里标记 COW，\code{P42.refcount = 2}。

\begin{longtable}{p{0.16\linewidth}p{0.34\linewidth}p{0.34\linewidth}}
\toprule
时刻 & 父进程 PTE & 子进程 PTE \\
\midrule
fork 前 & \code{0x7000 -> P42, writable} & 不存在 \\
fork 后 & \code{0x7000 -> P42, read-only, COW} & \code{0x7000 -> P42, read-only, COW} \\
子进程写入前 & \code{P42.refcount = 2} & store 触发写保护 fault \\
子进程 fault 后 & 仍指向 \code{P42, read-only, COW} & 指向新页 \code{P99, writable} \\
后续父进程写入 & 若 \code{P42.refcount = 1} 可直接恢复 writable & 已经和父进程分离 \\
\bottomrule
\end{longtable}

关键点是“写入必须在 store 提交前被拦住”。CPU 执行子进程的 \code{*p = 'x'} 时，地址翻译发现 PTE present 但不可写，于是保存 faulting address、错误码和 RIP 进入内核。因为 store 还没有修改 \code{P42}，内核才能安全复制旧页到 \code{P99}，把子进程 PTE 改为 \code{P99, writable}，再返回原指令重试。重试后写入落在 \code{P99}，父进程继续看到旧的 \code{"abc"}。

失败路径也要讲清楚。若 fault 地址不在允许写的 VMA 内，不能执行 COW，而应给进程发段错误。若分配 \code{P99} 失败，可能触发回收或 OOM。若更新 PTE 后没有失效旧 TLB，CPU 可能继续使用旧只读翻译，导致重复 fault；更危险的是撤销可写权限时没有失效旧可写翻译，会让某个 CPU 绕过新的只读保护继续写共享页。COW 的正确性不是“复制一页”这么简单，而是 VMA 权限、PTE 权限、物理页引用计数、TLB 失效和 fault 可重试性共同成立。

\subsection{页面置换}

当物理内存不足时，内核要回收页面。理想算法 OPT 总是淘汰未来最晚再访问的页面，但它需要知道未来，不能实现。LRU 淘汰最近最久未使用的页，更接近实际局部性，但精确 LRU 维护成本高。Clock 算法用访问位近似 LRU：扫描环形列表，访问位为 1 就清零并跳过，访问位为 0 就淘汰。

\begin{lstlisting}
clock_evict():
  while true:
    page = clock_hand.page
    if page.referenced:
      page.referenced = false
      advance(clock_hand)
    else if page.dirty:
      schedule_writeback(page)
      advance(clock_hand)
    else:
      remove_and_return(page)
\end{lstlisting}

置换策略必须和文件系统、匿名页、swap、dirty writeback 结合。干净文件页可以直接丢弃，未来再从文件读；脏文件页要先写回；匿名页若要回收，需要 swap 或压缩；被锁定或正在 I/O 的页不能随便回收。

\subsection{多级页表}

64 位地址空间太大，不能为每个虚拟页都预留一个线性页表。多级页表把虚拟页号拆成多段索引，只为实际使用的区域分配中间页表页。以四级页表为例，虚拟地址被拆成 L4、L3、L2、L1 和页内 offset。

\begin{lstlisting}
walk_page_table(root, vaddr, create):
  table = root
  for level in [L4, L3, L2]:
    index = index_for(level, vaddr)
    if table[index] not present:
      if not create:
        return null
      table[index] = allocate_page_table()
    table = table[index].next_table
  return &table[index_for(L1, vaddr)]
\end{lstlisting}

查找复杂度按层数是常数，但每层都可能带来内存访问。TLB 的价值就在于缓存最近翻译，避免每次访问都走多级页表。页表设计还影响内核内存占用：稀疏地址空间只分配用到的页表页，密集映射则消耗更多元数据。

\subsection{TLB shootdown}

改变页表后，CPU TLB 中可能仍缓存旧翻译。单核系统可以本地失效；多核系统中，同一地址空间可能正在其他 CPU 运行，需要发送 shootdown IPI 让它们失效对应条目。

\begin{lstlisting}
unmap_page(mm, vaddr):
  pte = walk_page_table(mm.root, vaddr, false)
  old = clear_pte(pte)
  cpus = mm.active_cpus
  for cpu in cpus except current_cpu:
    send_tlb_shootdown(cpu, mm, vaddr)
  local_invlpg(vaddr)
  wait_for_shootdown_ack(cpus)
  put_page(old.page)
\end{lstlisting}

顺序很重要。若在其他 CPU 失效 TLB 前就释放物理页，该页可能被重新分配给别的对象，而旧 CPU 仍通过过期 TLB 写入它，形成严重内存破坏。TLB shootdown 是虚拟内存和多核同步的交叉点。

再把这个危险展开成具体事故。进程 X 有虚拟页 \code{0x9000 -> P10}，它的两个线程分别在 CPU0 和 CPU1 上运行。CPU1 的 TLB 里已经缓存了这条翻译。CPU0 执行 \code{munmap(0x9000)}，把页表项清掉，然后若立刻把 \code{P10} 释放给物理页分配器，分配器可能马上把 \code{P10} 交给内核网络栈当作 packet buffer。此时 CPU1 如果还没收到 shootdown，仍可用旧 TLB 把用户态 \code{0x9000} 翻译到 \code{P10}，一次普通 store 就会写坏网络 buffer。这个 bug 表面可能表现成网络包校验失败、随机内核崩溃或安全越权，根因却是页表撤销和 TLB 旧副本之间的生命周期错误。

因此 shootdown 的完成点必须在“物理页可重用”之前。教学模型可以记成三步：先让新页表状态不可再产生新翻译；再让所有可能持有旧翻译的 CPU 丢掉旧 TLB；最后释放旧物理页或改变其用途。

\begin{longtable}{p{0.18\linewidth}p{0.38\linewidth}p{0.32\linewidth}}
\toprule
阶段 & 必须完成的事 & 如果提前进入下一步 \\
\midrule
清 PTE & 后续 page walk 不能再得到旧映射 & 新 TLB 项仍可能被硬件装入 \\
本地失效 & 当前 CPU 不再使用旧翻译 & 当前线程可能继续访问已撤销页 \\
远端 IPI & 其他 active CPU 执行 \code{invlpg} 或等价操作 & 远端 CPU 可能写入已释放物理页 \\
等待 ack & 确认旧翻译不可再被使用 & 释放页会造成 use-after-free 到物理内存 \\
释放旧页 & 页框回到分配器或引用计数下降 & 到这一步才安全改变页用途 \\
\bottomrule
\end{longtable}

这也是为什么 \code{mprotect}、\code{munmap}、COW 和进程退出在多核上不是“改一张表”就结束。内核还要知道哪些 CPU 可能正在运行这个地址空间，哪些映射范围需要 flush，能否批量合并 shootdown，是否允许延迟释放。优化可以减少 IPI 次数，但不能破坏上面的安全边界。

\subsection{overcommit 与 OOM}

虚拟内存允许系统承诺比物理内存更多的虚拟地址。例如 \code{mmap} 或 \code{malloc} 成功不代表物理页已经分配，首次写入才可能触发分配。overcommit 提高灵活性，但内存真正不足时必须选择失败策略。

\begin{lstlisting}
handle_anon_fault(vma, addr):
  page = alloc_page()
  if page == null:
    victim = oom_select_process()
    if victim == null:
      return SIGBUS_OR_ENOMEM
    kill(victim)
    retry allocation
  zero_page(page)
  install_pte(addr, page, WRITE | USER)
\end{lstlisting}

OOM killer 不是随意杀进程。它通常综合内存占用、权限、oom score、是否关键服务、cgroup 限制和可回收性。容器场景下，cgroup 内 OOM 应优先限制在该 cgroup 内，而不是扩大到整个宿主。

把 overcommit 放进一个常见服务事故。机器有 8GiB 内存，服务 A 启动时 \code{malloc(6GiB)} 做缓存索引，服务 B 也 \code{malloc(6GiB)} 准备批处理。两个 \code{malloc} 都可能成功，因为内核只是给它们各自建立虚拟区间和 VMA，并没有立即分配 12GiB 物理页。监控里 VSZ 很大，但 RSS 还不高。事故发生在两者开始真正写入页面时：

\begin{lstlisting}
for each 4096-byte page in allocated range:
  p[i] = 0;
\end{lstlisting}

每写一个尚未物化的匿名页，CPU 都会触发缺页；内核确认 VMA 允许写，分配物理页，清零，安装 PTE，然后重试 store。刚开始这很顺利；随着空闲页减少，内核先回收干净文件页，再尝试写回脏页，再考虑 swap 或压缩；若仍不能满足分配，就进入 OOM 决策。此时失败点不是 \code{malloc}，而是某一次普通 store 的缺页处理。

\begin{longtable}{p{0.18\linewidth}p{0.38\linewidth}p{0.32\linewidth}}
\toprule
阶段 & 内核状态 & 应用看到什么 \\
\midrule
\code{malloc} 成功 & VMA/堆区扩展，物理页可能未分配 & 返回非空指针 \\
首次写前几页 & 匿名页缺页，分配物理页并清零 & store 正常完成，只是可能变慢 \\
内存压力上升 & 回收 page cache、写回脏页、触发 swap & 延迟抖动，RSS 上升 \\
无法回收 & OOM 在进程或 cgroup 内选择受害者 & 某进程被杀，或 fault 返回失败语义 \\
cgroup 限制命中 & 只在该资源组内回收和 OOM & 宿主其他服务不应被拖下水 \\
\bottomrule
\end{longtable}

这解释了为什么服务不能只用“malloc 成功”作为容量保证。真正可靠的设计要设置 cgroup memory.max、应用层缓存上限、分批预热、失败回退和监控 RSS/PSI。overcommit 的价值是让稀疏使用的大地址空间可行；代价是失败可能延迟到第一次触碰具体页面时才出现。

\section{内存压力、回收失败与 OOM}

测试要覆盖成功翻译、写只读页失败、访问未映射页失败、unmap 后访问失败、页内 offset 保留。错误结果必须携带可读诊断，不能只返回一个空值。

\begin{itemize}
\tightlist
\item \code{fork} 后父子共享只读 COW 页，任一方写入后得到私有副本。
\item \code{mmap} 文件页首次访问触发填页，第二次访问走热路径。
\item 栈增长必须受上限约束，不能无限扩张。
\item \code{mprotect} 改变权限后，旧访问策略失效。
\item \code{munmap} 后再次访问必须失败。
\item 用户态指针传入系统调用时，内核必须检查可读/可写范围。
\end{itemize}

权限测试尤其重要。没有权限测试的虚拟内存模型，会把地址转换误写成普通哈希表查找，丢掉操作系统最重要的保护语义。

\subsection{地址策略、当前映射与页面归属}

虚拟内存的细节很多：page fault、VMA 切分、PTE 标志、TLB shootdown、COW、swap、THP、mlock、PSI、OOM 和 cgroup。主线里先把它们压成一张“缺页账本”。每次 CPU 访问地址失败，内核都要回答：地址属于哪个 VMA，访问类型是否被允许，页面内容从哪里来，PTE 怎样安装，旧 TLB 是否要失效，失败时返回信号还是 errno，线程是否需要睡眠等待 I/O 或回收。

\begin{longtable}{p{0.20\linewidth}p{0.38\linewidth}p{0.30\linewidth}}
\toprule
场景 & 内核主要动作 & 关键风险 \\
\midrule
匿名首次写 & 分配物理页、清零、安装 writable PTE & overcommit 后在首次触碰时 OOM \\
COW 写 fault & 检查引用、复制页面、替换 PTE、flush TLB & 父子进程互相污染或复制后权限错误 \\
文件映射 fault & 查 page cache、必要时发起读 I/O、安装映射 & 把 I/O 等待误判成 CPU 慢 \\
\code{mprotect} & VMA 切分、改权限、更新页表、失效 TLB & 旧 TLB 保留过宽权限 \\
内存回收 & 扫 LRU、写回脏页、swap out、唤醒等待者 & 回收不可回收页或引发严重抖动 \\
\bottomrule
\end{longtable}

虚拟内存可以按对象追踪：\code{mm\_struct} 是进程地址空间账本，VMA 是区间权限和来源
账本，PTE 是处理器当前可执行的翻译账本，\code{struct page/folio} 是物理页生命周期账本，
rmap 则把物理页反查到映射它的地址空间。任何优化都不能破坏这些账本的一致性。THP 减少
TLB 压力，但拆分时仍要恢复普通页语义；swap 降低内存占用，但 fault 延迟会进入请求尾部；
mlock 保护实时路径或密钥页，但会减少可回收集合；PSI 不消除内存压力，只报告线程因内存
等待而失去前进能力。

内核自己的内存分配也属于这张账本。页分配器负责按页框分配和回收，buddy 系统处理连续页块；slab/SLUB 把页切成固定大小对象，服务 \code{task\_struct}、inode、dentry、socket 等高频对象；\code{vmalloc} 提供虚拟连续但物理不连续的内核映射。教学实现可以先只有 page allocator 和几个对象池，但必须保留三个边界：中断上下文不能随意睡眠等待内存，DMA buffer 可能要求物理连续或满足设备地址限制，对象缓存释放后不能仍被 RCU 读者或异步 completion 引用。

\begin{longtable}{p{0.20\linewidth}p{0.38\linewidth}p{0.30\linewidth}}
\toprule
分配器 & 适合的对象 & 验收问题 \\
\midrule
页分配器 & 页表页、大块缓冲、匿名页和 page cache 页 & 压力下是否能回收、写回或明确失败 \\
SLUB/slab & 大量同类型小对象 & 析构、引用计数、越界和 use-after-free 是否可检测 \\
\code{vmalloc} & 大型内核虚拟连续缓冲 & 不能交给要求物理连续 DMA 的设备 \\
per-CPU allocator & 每 CPU 统计、队列和热路径缓存 & CPU 热插拔和迁移时账本是否一致 \\
\bottomrule
\end{longtable}

\begin{rubricbox}
本章验收问题：给一个虚拟地址和一次读/写/取指访问，能否说明它命中哪个 VMA、PTE 应有哪组权限、缺页后页面从哪里来、是否可能睡眠、失败给用户什么信号或错误、哪些观测指标能证明判断。能回答这些，才算把虚拟内存读成 OS 机制，而不是地址转换表。
\end{rubricbox}

\section{巨页、交换与地址翻译成本}

\subsection{页表项标志}

x86-64 PTE 典型标志包括 present、writable、user、accessed、dirty、global、NX 等。操作系统把 VMA 权限、文件打开模式、COW 状态和硬件标志合成 PTE。注意 COW 不是单纯硬件位，通常由软件在只读 PTE 和附加元数据中表达。

\begin{longtable}{p{0.22\linewidth}p{0.48\linewidth}}
\toprule
标志 & 语义 \\
\midrule
present & PTE 可用于地址翻译，否则触发缺页 \\
writable & store 是否允许 \\
user & ring 3 是否可访问 \\
accessed & CPU 访问后置位，用于回收近似 LRU \\
dirty & CPU 写入后置位，用于写回和 COW 判断 \\
NX & 禁止取指执行 \\
\bottomrule
\end{longtable}

权限收缩必须刷新 TLB。例如把页从可写改成只读，如果旧 TLB 仍保留可写权限，写入可能绕过新的 COW 或 mprotect 规则。

\subsection{mprotect、munmap 和 VMA 切分}

\code{mprotect} 改变区间权限，\code{munmap} 删除映射。二者都可能把一个 VMA 切成三段：前半保持原权限，中间修改或删除，后半保持原权限。

\begin{lstlisting}
before:
  [0x1000, 0x9000) RW
mprotect [0x3000, 0x5000) R
after:
  [0x1000, 0x3000) RW
  [0x3000, 0x5000) R
  [0x5000, 0x9000) RW
\end{lstlisting}

VMA 切分后，要更新页表权限、失效 TLB、处理锁定页和文件映射引用。教学实现若只改一个 flags 字段，会在部分区间修改时出错。

\subsection{Transparent Huge Page}

普通页常为 4KiB，大页可为 2MiB 或更大。THP 尝试把连续匿名页合并成 PMD 级大页，减少 TLB miss。代价是分配连续内存更难，拆分和 COW 成本更高，延迟可能出现尖峰。

\begin{lstlisting}
khugepaged:
  scan VMAs marked hugepage-friendly
  find 512 contiguous 4KiB pages
  allocate 2MiB huge page
  copy or remap contents
  replace PTE table with PMD huge mapping
\end{lstlisting}

THP 是性能优化，不应改变用户语义。发生 mprotect、munmap、COW 或内存压力时，大页可能被拆回普通页。性能报告要说明 THP 是否开启，否则 TLB 和内存延迟数据难以比较。

\subsection{swap cache 和 swappiness}

匿名页回收需要把内容写到 swap。swap slot 是匿名页在交换区的位置；swap cache 避免同一 swap 页被重复读写。swappiness 影响内核在文件页和匿名页之间的回收倾向。

\begin{lstlisting}
swap_out(page):
  slot = allocate_swap_slot()
  write_page_to_swap(slot, page)
  replace_pte_with_swap_entry(vaddr, slot)
  free_physical_page(page)

swap_in(fault):
  slot = fault.swap_entry
  page = swap_cache_lookup_or_read(slot)
  install_pte(fault.addr, page)
\end{lstlisting}

swap 让系统在内存压力下继续运行，但延迟可能变得非常高。服务端系统常用 PSI、cgroup memory.high 和限流来避免进入严重 swap 抖动。

\subsection{mlock 和不可回收页}

\code{mlock} 把页锁在内存中，常用于实时路径、密码材料或避免 page fault 的场景。锁定页不能被普通回收换出，因此必须受权限和 rlimit 限制。

\begin{lstlisting}
mlock(addr, len):
  check RLIMIT_MEMLOCK
  for each VMA page in range:
    fault page in if needed
    mark unevictable
\end{lstlisting}

滥用 mlock 会减少可回收内存，影响整机。教学 OS 要把”不可回收”作为页面状态加入回收测试，证明回收器不会错误丢弃 locked page。

\inputdetail{chapters/ch11-memory-syscall-deep-dive}
